\documentclass[12pt]{article}
\usepackage[utf8]{inputenc}
\usepackage[spanish]{babel}
\usepackage{amsmath}
\usepackage{amssymb}
\usepackage{booktabs}
\usepackage{geometry}
\geometry{margin=2.5cm}
\usepackage[most]{tcolorbox}
\newtcolorbox{intuicion}{
  colback=blue!5!white,
  colframe=blue!40!black,
  fonttitle=\bfseries,
  title=Intuici\'on,
  boxrule=0.6pt,
  arc=4pt,
  left=6pt, right=6pt, top=4pt, bottom=4pt
}
\newtcolorbox{intuicionmat}{
  colback=green!5!white,
  colframe=green!40!black,
  fonttitle=\bfseries,
  title=\textit{?`Qu\'e dice esta ecuaci\'on?},
  boxrule=0.6pt,
  arc=4pt,
  left=6pt, right=6pt, top=4pt, bottom=4pt
}
\newtcolorbox{explicacion}{
  colback=orange!5!white,
  colframe=orange!50!black,
  fonttitle=\bfseries,
  title=Explicaci\'on,
  boxrule=0.6pt,
  arc=4pt,
  left=6pt, right=6pt, top=4pt, bottom=4pt
}

\title{\textbf{Resumen: Cap\'itulo 4 --- El Teorema de Arbitraje}\\[4pt]
\large \textit{An Undergraduate Introduction to Financial Mathematics}\\
\normalsize J.~R.~Buchanan}
\author{}
\date{Junio 2026}

\begin{document}
\maketitle
\tableofcontents
\newpage

%%% =========================================================
\section{Introducci\'on: El Concepto de Arbitraje (The Concept of Arbitrage)}
%%% =========================================================

El \textbf{arbitraje (arbitrage)} es la situaci\'on que surge cuando dos
instrumentos financieros est\'an incorrectamente valuados entre s\'i, haciendo
posible obtener una ganancia financiera sin riesgo. Por ejemplo, si el banco A
ofrece pr\'estamos al 5\,\% y el banco B ofrece dep\'ositos al 6\,\%, es posible
tomar un pr\'estamo del banco A, depositarlo en el banco B y obtener el 1\,\%
de diferencia como beneficio puro.

Los mercados financieros eficientes impiden que estas oportunidades sean
frecuentes: cuando surge el arbitraje, los inversores acuden masivamente a
los instrumentos con precio incorrecto y el mercado corrige los precios.
El \textbf{Teorema de Arbitraje (Arbitrage Theorem)} establece que o bien
las probabilidades de los resultados son tales que todas las apuestas son
justas, o bien existe un esquema de apuestas que produce una ganancia positiva
independiente del resultado. Este resultado es la base de la derivaci\'on de
la ecuaci\'on de Black--Scholes que se estudia en el Cap\'itulo~6.

%%% =========================================================
\section{Momios y Probabilidades (Odds and Probabilities)}
%%% =========================================================

\subsection{Relaci\'on entre Momios y Probabilidades}

Para simplificar el c\'alculo de pagos se utilizan \textbf{momios (odds)}
en lugar de probabilidades directas. Si los momios en contra de un resultado
son $n:1$, la probabilidad de ese resultado es:

\begin{equation}
  P = \frac{1}{n+1}
\end{equation}

\begin{intuicionmat}
Si los momios son ``2 a 1 en contra'', de cada 3 escenarios posibles el
resultado ocurre 1 vez, de modo que su probabilidad es $1/3$. La f\'ormula
simplemente cuenta cu\'antas veces aparece el resultado favorable sobre el
total de escenarios impl\'icitos en los momios.
\end{intuicionmat}

De forma m\'as general, si los momios son $n:m$ (el evento no ocurre $n$
veces y ocurre $m$ veces), la probabilidad del evento es:

\begin{equation}
  P = \frac{m}{n+m}
\end{equation}

\begin{intuicionmat}
Los momios $n:m$ son equivalentes a los momios $\frac{n}{m}:1$. La
probabilidad es la fracci\'on de veces que ocurre el evento sobre el total
de resultados impl\'icitos. Esta generalizaci\'on permite trabajar con
cualquier par de enteros naturales $n$ y $m$.
\end{intuicionmat}

El mecanismo de pago es el siguiente: si los momios en contra de un resultado
son $n:1$ y se realiza una apuesta unitaria, se ganan $n$ unidades cuando el
resultado ocurre y se pierde la unidad apostada en caso contrario. Los pagos
escalan multiplicativamente para apuestas no unitarias.

\begin{intuicion}
Los momios son simplemente una forma pr\'actica de expresar la rareza de un
evento: cuanto mayores son los momios en contra, m\'as improbable es el
resultado y mayor es la ganancia potencial al apostar por \'el. Si las
probabilidades impl\'icitas en los momios de todos los resultados no suman
exactamente 1, las tasas est\'an en desequilibrio y existe una oportunidad
de arbitraje --- una estrategia que gana sin importar qu\'e resultado ocurra.
\end{intuicion}

\subsection{Ejemplo: Partida de Tenis y Arbitraje}

Suponga que los momios en contra de que el jugador A derrote al jugador B
son $\frac{3}{2}:1$ y los momios en contra de que B derrote a A son
$\frac{3}{7}:1$. Convirtiendo a probabilidades:

\begin{align}
  P(A \text{ gana}) &= \frac{1}{\frac{3}{2}+1} = 0{,}4 \\[4pt]
  P(B \text{ gana}) &= \frac{1}{\frac{3}{7}+1} = 0{,}7
\end{align}

\begin{intuicionmat}
Las dos probabilidades suman $0{,}4 + 0{,}7 = 1{,}1 \neq 1$, lo cual es
incoherente: las probabilidades de todos los resultados de un experimento
deben sumar exactamente 1. Esta inconsistencia es la se\~nal de que existen
momios incorrectos y por lo tanto existe arbitraje.
\end{intuicionmat}

Apostando $-1$ al jugador~A (es decir, apostando $1$ en contra de A) y
$-2$ al jugador~B, se obtiene:
\begin{itemize}
  \item Si gana A: se pierde $3/2$ en la primera apuesta y se gana 2 en la
        segunda, para una ganancia neta de $1/2$.
  \item Si gana B: se gana 1 en la primera apuesta y se pierde $6/7$ en la
        segunda, para una ganancia neta de $1/7$.
\end{itemize}

En ambos casos la ganancia es positiva, lo que confirma la existencia de
arbitraje.

\begin{intuicion}
Este ejemplo ilustra el Teorema de Arbitraje en su forma m\'as concreta:
cuando las probabilidades impl\'icitas en los momios no son coherentes, siempre
existe una combinaci\'on de apuestas que garantiza ganancia. Encontrar esa
combinaci\'on en situaciones complejas requiere las herramientas de
programaci\'on lineal desarrolladas en la siguiente secci\'on.
\end{intuicion}

%%% =========================================================
\section{Teorema de Dualidad de la Programaci\'on Lineal\\
         (Duality Theorem of Linear Programming)}
%%% =========================================================

\subsection{Formulaci\'on del Problema de Programaci\'on Lineal}

Un \textbf{problema cl\'asico de programaci\'on lineal (classic linear
programming problem)} consiste en minimizar una funci\'on de costo lineal
sujeta a que la soluci\'on sea \textbf{factible (feasible)}.

El vector columna $\mathbf{x} \in \mathbb{R}^{n \times 1}$ es factible si
$x_i \geq 0$ para $i = 1, 2, \ldots, n$ y satisface las ecuaciones de
restricci\'on $A\mathbf{x} = \mathbf{b}$, donde $A$ es una matriz $m \times n$
y $\mathbf{b} \in \mathbb{R}^{m \times 1}$. El vector fila
$\mathbf{c} \in \mathbb{R}^{1 \times n}$ define la \textbf{funci\'on de costo
(cost function)}:

\begin{equation}
  \mathbf{c} \cdot \mathbf{x} = c_1 x_1 + c_2 x_2 + \cdots + c_n x_n
\end{equation}

\begin{intuicionmat}
Este producto punto es el ``costo total'' de una combinaci\'on de recursos:
$c_i$ es el costo unitario del recurso $i$ y $x_i$ es la cantidad usada.
La programaci\'on lineal busca la combinaci\'on de cantidades m\'as barata
que cumpla simult\'aneamente todas las restricciones.
\end{intuicionmat}

Las \textbf{restricciones de desigualdad (inequality constraints)} se pueden
convertir en restricciones de igualdad introduciendo \textbf{variables de
holgura (slack variables)}. Por ejemplo:

\begin{equation}
  x_1 + x_2 + x_3 \leq 1
  \quad \xrightarrow{\text{variable de holgura } x_4 \geq 0} \quad
  x_1 + x_2 + x_3 + x_4 = 1
\end{equation}

\begin{intuicionmat}
La variable de holgura $x_4 \geq 0$ absorbe la diferencia entre el lado
izquierdo y el derecho de la desigualdad, convirtiendo el problema en uno
de igualdad sin alterar el conjunto de soluciones factibles. No es necesario
modificar la funci\'on de costo (el coeficiente de $x_4$ es cero).
\end{intuicionmat}

El conjunto de puntos factibles siempre forma un \textbf{conjunto convexo
(convex set)}: un conjunto donde el segmento de l\'inea que une cualquier
par de puntos del conjunto tambi\'en pertenece al conjunto. La soluci\'on
\'optima se encuentra en un v\'ertice (esquina) de este conjunto convexo.

\begin{intuicion}
La programaci\'on lineal busca el punto m\'as barato dentro de una regi\'on
geom\'etrica definida por las restricciones. Como la funci\'on de costo es
lineal y el conjunto factible es convexo, la soluci\'on \'optima siempre
est\'a en un v\'ertice --- nunca en el interior. Esto hace que el problema
sea tratable incluso con miles de variables, pues basta examinar los v\'ertices.
\end{intuicion}

\subsection{El Par Primal--Dual (The Primal--Dual Pair)}

Para cada \textbf{problema primal (primal problem)} de programaci\'on lineal
existe un \textbf{problema dual (dual problem)} asociado:

\medskip
\noindent\textbf{Primal:} Minimizar $\mathbf{c} \cdot \mathbf{x}$ sujeto a
$A\mathbf{x} = \mathbf{b}$ y $\mathbf{x} \geq 0$.\\
\textbf{Dual:} Maximizar $\mathbf{y} \cdot \mathbf{b}$ sujeto a
$\mathbf{y}A \leq \mathbf{c}$.
\medskip

Las relaciones estructurales entre ambos problemas son:
\begin{enumerate}
  \item El desconocido del dual es el vector fila $\mathbf{y} \in \mathbb{R}^{1 \times n}$.
  \item El vector $\mathbf{b}$ pasa de la restricci\'on del primal a la
        funci\'on objetivo del dual.
  \item El vector $\mathbf{c}$ pasa de la funci\'on de costo del primal a
        la restricci\'on del dual.
  \item Las restricciones del dual son desigualdades; hay $n$ de ellas.
  \item No existe restricci\'on de no negatividad ($\mathbf{y} \geq 0$) en el
        dual est\'andar.
\end{enumerate}

\begin{intuicion}
El dual es un ``espejo'' del primal. En econom\'ia, las variables duales tienen
interpretaci\'on como \textbf{precios sombra (shadow prices)}: el valor marginal
de relajar cada restricci\'on del primal en una unidad. Resolver el dual suele
ser computacionalmente m\'as eficiente cuando tiene menos restricciones que el
primal. La simetr\'ia profunda entre ambos problemas es lo que hace posible
el Teorema de Arbitraje.
\end{intuicion}

\subsection{Teorema de Dualidad D\'ebil (Weak Duality Theorem)}

\textbf{Teorema 4.1 (Dualidad D\'ebil / Weak Duality Theorem):} Si $\mathbf{x}$
e $\mathbf{y}$ son los desconocidos del primal y del dual respectivamente,
entonces:

\begin{equation}
  \mathbf{y} \cdot \mathbf{b} \leq \mathbf{c} \cdot \mathbf{x}
\end{equation}

Si $\mathbf{y} \cdot \mathbf{b} = \mathbf{c} \cdot \mathbf{x}$, entonces
estos vectores son \'optimos para sus respectivos problemas.

\begin{intuicionmat}
El valor objetivo del dual es siempre una \textbf{cota inferior (lower bound)}
del valor objetivo del primal. Esto significa que cualquier soluci\'on factible
del dual nos da un l\'imite garantizado para el m\'inimo del primal. Cuando
la cota alcanza el valor del primal --- es decir, cuando hay igualdad ---
hemos encontrado el \'optimo global de ambos problemas simult\'aneamente.
\end{intuicionmat}

\textit{Demostraci\'on.} Las soluciones deben satisfacer $A\mathbf{x} = \mathbf{b}$,
$x_i \geq 0$ y $\mathbf{y}A \leq \mathbf{c}$. Multiplicando la restricci\'on
del primal por $\mathbf{y}$ y la del dual por $\mathbf{x}$:

\begin{equation}
  \mathbf{y}A\mathbf{x} = \mathbf{y} \cdot \mathbf{b}
  \qquad \text{y} \qquad
  \mathbf{y}A\mathbf{x} \leq \mathbf{c} \cdot \mathbf{x}
\end{equation}

\begin{intuicionmat}
La expresi\'on $\mathbf{y}A\mathbf{x}$ es el v\'inculo entre primal y dual:
es simult\'aneamente igual a $\mathbf{y} \cdot \mathbf{b}$ (por la restricci\'on
del primal) y menor o igual que $\mathbf{c} \cdot \mathbf{x}$ (por la restricci\'on
del dual y la no negatividad de $\mathbf{x}$). Encadenar estas dos relaciones
da $\mathbf{y} \cdot \mathbf{b} \leq \mathbf{c} \cdot \mathbf{x}$.
\end{intuicionmat}

Concatenando: $\mathbf{y} \cdot \mathbf{b} = \mathbf{y}A\mathbf{x} \leq \mathbf{c} \cdot \mathbf{x}$.
Si $\mathbf{y} \cdot \mathbf{b} = \mathbf{c} \cdot \mathbf{x}$, ning\'un
$\mathbf{y}$ puede hacer $\mathbf{y} \cdot \mathbf{b}$ mayor ni ning\'un
$\mathbf{x}$ puede hacer $\mathbf{c} \cdot \mathbf{x}$ menor, por lo que
ambos son \'optimos. $\square$

\begin{intuicion}
La elegancia de esta demostraci\'on reside en que un \'unico t\'ermino
$\mathbf{y}A\mathbf{x}$ conecta los dos problemas. La no negatividad de
$\mathbf{x}$ es esencial: sin ella, la desigualdad $\mathbf{y}A \leq \mathbf{c}$
al multiplicar por $\mathbf{x}$ no preservar\'ia el sentido. Esta hip\'otesis
es lo que distingue el primal est\'andar del dual, que no exige $\mathbf{y} \geq 0$.
\end{intuicion}

\subsection{Condici\'on de Optimalidad y Holgura Complementaria\\
            (Optimality Condition and Complementary Slackness)}

Como consecuencia del Teorema de Dualidad D\'ebil, cuando $\mathbf{x}$ e
$\mathbf{y}$ son \'optimos se cumple:

\begin{align}
  \mathbf{y} \cdot \mathbf{b} &= \mathbf{y}A\mathbf{x} = \mathbf{c} \cdot \mathbf{x} \nonumber \\
  \Rightarrow\quad (\mathbf{c} - \mathbf{y}A) \cdot \mathbf{x} &= 0
  \label{eq:holgura}
\end{align}

\begin{intuicionmat}
El vector $\mathbf{c} - \mathbf{y}A$ mide la ``holgura de costo'': cu\'anto
falta para que la restricci\'on dual sea activa en cada componente. La
ecuaci\'on dice que el producto punto de esta holgura con la soluci\'on
$\mathbf{x}$ es exactamente cero. Como ambos vectores son no negativos,
esto s\'olo puede ocurrir si, para cada componente $j$, al menos uno de
los dos factores es cero.
\end{intuicionmat}

\textbf{Teorema 4.2 (Condici\'on de Optimalidad / Optimality Condition):}
La optimalidad en los problemas primal y dual requiere que para cada
$j = 1, \ldots, n$:

\begin{equation}
  x_j = 0 \qquad \text{o} \qquad (\mathbf{y}A)_j = c_j
\end{equation}

\begin{intuicionmat}
Este es el teorema de \textbf{holgura complementaria (complementary slackness)}:
o bien la variable de decisi\'on $x_j$ es cero (no se usa ese recurso), o
bien la restricci\'on dual $j$-\'esima es activa (el costo marginal iguala el
precio). Una variable positiva con restricci\'on dual inactiva es imposible
en el \'optimo, y esto permite reducir dr\'asticamente el primal una vez
conocida la soluci\'on dual.
\end{intuicionmat}

\begin{intuicion}
Este teorema es una herramienta pr\'actica poderosa. En el Ejemplo~4.2 del
texto, conocer el \'optimo del dual $(y_1, y_2) = (2,3)$ revela que las
restricciones 1 y 3 del dual son estrictamente inactivas. Por la condici\'on
de optimalidad, $x_1 = x_3 = 0$ en el primal. El problema de 4 variables se
reduce a uno de 2 variables, resoluble por inspecci\'on: $(x_1, x_2, x_3, x_4)
= (0, 5, 0, 3)$ con valor \'optimo 19.
\end{intuicion}

\subsection{Teorema de Dualidad (Duality Theorem)}

\textbf{Teorema 4.3 (Dualidad / Duality Theorem) [Gale \textit{et al.}, 1951]:}
Si existe un $\mathbf{x}$ \'optimo en el primal, entonces existe un $\mathbf{y}$
\'optimo en el dual y:

\begin{equation}
  \min\,(\mathbf{c} \cdot \mathbf{x}) = \max\,(\mathbf{y} \cdot \mathbf{b})
\end{equation}

\begin{intuicionmat}
Esta es la versi\'on fuerte de la dualidad: el m\'inimo del primal coincide
exactamente con el m\'aximo del dual. No es una cota aproximada --- es igualdad
exacta. Esta propiedad de ``brecha cero'' entre primal y dual (\textit{strong
duality}) es lo que hace posible la demostraci\'on del Teorema de Arbitraje:
la factibilidad del primal equivale a la existencia de probabilidades justas,
y la ausencia de factibilidad equivale a la existencia de arbitraje.
\end{intuicionmat}

El teorema admite una formulaci\'on sim\'etrica: si cualquiera de los dos
problemas tiene soluci\'on \'optima, el otro tambi\'en la tiene y sus valores
\'optimos son iguales ($\mathbf{c} \cdot \mathbf{x} = \mathbf{y} \cdot \mathbf{b}$).
Cuando el primal tiene \textbf{restricciones de desigualdad (inequality
constraints)} $A\mathbf{x} \geq \mathbf{b}$, introducir variables de holgura
produce la \textbf{forma sim\'etrica (symmetric form)}:

\medskip
\noindent\textbf{Primal sim\'etrico:} Minimizar $\mathbf{c} \cdot \mathbf{x}$
s.a. $A\mathbf{x} \geq \mathbf{b}$, $\mathbf{x} \geq 0$.\\
\textbf{Dual sim\'etrico:} Maximizar $\mathbf{y} \cdot \mathbf{b}$ s.a.
$\mathbf{y}A \leq \mathbf{c}$, $\mathbf{y} \geq 0$.
\medskip

\begin{intuicion}
En la forma sim\'etrica ambos problemas tienen restricciones de desigualdad
y no negatividad de sus variables. Esta presentaci\'on hace evidente que el
dual del dual es el primal --- el par es completamente reversible. Adem\'as,
en la forma sim\'etrica la condici\'on $\mathbf{y} \geq 0$ permite interpretar
$\mathbf{y}$ directamente como un vector de probabilidades, lo que es la clave
de la conexi\'on con el Teorema de Arbitraje.
\end{intuicion}

%%% =========================================================
\section{El Teorema Fundamental de las Finanzas\\
         (The Fundamental Theorem of Finance)}
%%% =========================================================

\subsection{Configuraci\'on del Experimento de Apuestas}

Considere un experimento con $m$ posibles resultados numerados del $1$ al $m$.
Se pueden realizar $n$ apuestas (numeradas del $1$ al $n$) sobre los resultados.
Sea $r_{ij}$ el retorno por una apuesta unitaria en la apuesta $i$ cuando el
resultado del experimento es $j$.

El vector $\mathbf{x} = (x_1, x_2, \ldots, x_n)$ se denomina
\textbf{estrategia de apuesta (betting strategy)}: $x_i$ es el monto apostado
en la apuesta $i$. El retorno compuesto bajo el resultado $j$ es:

\begin{equation}
  R_j = \sum_{i=1}^{n} x_i r_{ij}
\end{equation}

\begin{intuicionmat}
Esta ecuaci\'on calcula la ganancia total bajo el resultado $j$ como la
suma ponderada de los retornos individuales. Es una combinaci\'on lineal:
si apostamos $x_i$ unidades en la apuesta $i$ y esa apuesta paga $r_{ij}$
por unidad cuando ocurre el resultado $j$, la ganancia total es la suma de
todos esos pagos. La estrategia de arbitraje busca que $R_j > 0$ para
todo $j$ simult\'aneamente.
\end{intuicionmat}

\begin{intuicion}
Una estrategia de apuesta es an\'aloga a un portafolio financiero: $x_i$
es la posici\'on en el activo $i$ y $r_{ij}$ es su retorno bajo el escenario $j$.
El Teorema de Arbitraje clasifica todos los sistemas posibles de precios en
dos categor\'ias mutuamente excluyentes: los que admiten una medida de
probabilidad ``justa'' y los que permiten una ganancia garantizada sin riesgo.
Esta dicotom\'ia es el fundamento de toda la teor\'ia moderna de valoraci\'on
de activos.
\end{intuicion}

\subsection{El Teorema de Arbitraje (The Arbitrage Theorem)}

\textbf{Teorema 4.4 (Teorema de Arbitraje / Arbitrage Theorem):} Exactamente
una de las dos afirmaciones siguientes es verdadera:

\begin{enumerate}
  \item Existe un vector de probabilidades $\mathbf{y} = (y_1, y_2, \ldots, y_m)$
        con $y_j \geq 0$ y $\sum_{j=1}^{m} y_j = 1$, tal que para toda apuesta
        $i = 1, 2, \ldots, n$:
        \begin{equation}
          \sum_{j=1}^{m} y_j \, r_{ij} = 0
          \label{eq:prob_justas}
        \end{equation}

  \item Existe una estrategia de apuesta $\mathbf{x} = (x_1, x_2, \ldots, x_n)$
        tal que para todo resultado $j = 1, 2, \ldots, m$:
        \begin{equation}
          \sum_{i=1}^{n} x_i \, r_{ij} > 0
          \label{eq:arbitraje}
        \end{equation}
\end{enumerate}

\begin{intuicionmat}
La ecuaci\'on~\eqref{eq:prob_justas} dice que existe una distribuci\'on de
probabilidad sobre los resultados tal que el valor esperado de cada apuesta
es exactamente cero: todos los precios son ``justos'' bajo esa medida. La
ecuaci\'on~\eqref{eq:arbitraje} dice que existe una combinaci\'on de apuestas
con ganancia positiva garantizada bajo cualquier resultado posible. El teorema
establece que estos dos casos se excluyen mutuamente y son exhaustivos: no
puede haber precios justos y arbitraje simult\'aneamente, ni puede ocurrir
que no exista ninguno de los dos.
\end{intuicionmat}

\begin{intuicion}
El Teorema de Arbitraje es la versi\'on financiera del Teorema de Dualidad.
La existencia de probabilidades justas corresponde a la factibilidad del
primal con valor cero; la existencia de una estrategia de arbitraje corresponde
a que el primal no es factible y el dual puede alcanzar cualquier valor positivo.
En las finanzas modernas, la condici\'on~\eqref{eq:prob_justas} define la
\textbf{medida de probabilidad neutral al riesgo (risk-neutral probability
measure)}, que es la herramienta central de la valoraci\'on de derivados.
\end{intuicion}

\subsection{Demostraci\'on del Teorema de Arbitraje}

La demostraci\'on reduce el Teorema de Arbitraje al Teorema de Dualidad
(Teorema~4.3). Se augmenta el vector de estrategia con un elemento $x_{n+1}$
que representa la ganancia positiva m\'inima garantizada a maximizar, sujeto a:

\begin{equation}
  \sum_{i=1}^{n} x_i r_{ij} \geq x_{n+1}, \qquad j = 1, 2, \ldots, m
\end{equation}

\begin{intuicionmat}
Al introducir $x_{n+1}$ como la ganancia m\'inima garantizada, el problema
de encontrar una estrategia de arbitraje (ganancia positiva para todo resultado)
se convierte en maximizar $x_{n+1}$. Si el m\'aximo es positivo, hay arbitraje;
si el m\'aximo es cero, no lo hay. Esta reformulaci\'on permite aplicar
directamente el Teorema de Dualidad.
\end{intuicionmat}

La maximizaci\'on de $x_{n+1}$ se escribe como: Maximizar $\mathbf{x} \cdot \mathbf{b}$
sujeto a $\mathbf{x}A \leq \mathbf{c}$, donde $\mathbf{c} = (0,\ldots,0)$,
$\mathbf{b} = (0,\ldots,0,1)$ y la matriz $A$ tiene la forma:

\begin{equation}
  A = \begin{bmatrix}
    -r_{11} & -r_{12} & \cdots & -r_{1m} \\
    -r_{21} & -r_{22} & \cdots & -r_{2m} \\
    \vdots  & \vdots  &        & \vdots  \\
    -r_{n1} & -r_{n2} & \cdots & -r_{nm} \\
    1       & 1       & \cdots & 1
  \end{bmatrix}
  \label{eq:matrizA}
\end{equation}

\begin{intuicionmat}
Las primeras $n$ filas de $A$ codifican los retornos negativados de cada
apuesta bajo cada resultado. La \'ultima fila de unos implementa la restricci\'on
$\sum_j x_{n+1} \leq m$ (una unidad por cada resultado), que con $\mathbf{c} = \mathbf{0}$
es la \'unica restricci\'on no trivial sobre $x_{n+1}$. La estructura en
bloques refleja la dualidad entre apuestas (filas) y resultados (columnas).
\end{intuicionmat}

Este es el dual del problema primal: Minimizar $\mathbf{c} \cdot \mathbf{y} = 0$
sujeto a $A\mathbf{y} = \mathbf{b}$ y $\mathbf{y} \geq 0$, que en forma
matricial es:

\begin{equation}
  \begin{bmatrix}
    -r_{11} & \cdots & -r_{n1} & 1 \\
    -r_{12} & \cdots & -r_{n2} & 1 \\
    \vdots  &        & \vdots  & \vdots \\
    -r_{1m} & \cdots & -r_{nm} & 1
  \end{bmatrix}
  \begin{bmatrix} y_1 \\ y_2 \\ \vdots \\ y_{m-1} \\ y_m \end{bmatrix}
  =
  \begin{bmatrix} 0 \\ 0 \\ \vdots \\ 0 \\ 1 \end{bmatrix}
\end{equation}

\begin{intuicionmat}
Este sistema de ecuaciones es exactamente la condici\'on de que $\mathbf{y}$
sea un vector de probabilidades ($y_j \geq 0$, $\sum_j y_j = 1$) con valor
esperado cero para todas las apuestas. Las primeras $m$ ecuaciones dicen
$\sum_{i} y_i (-r_{ji}) = 0$, es decir $\sum_i y_i r_{ji} = 0$; la \'ultima
ecuaci\'on dice $\sum_i y_i = 1$.
\end{intuicionmat}

El primal es factible con valor m\'inimo cero si y solo si $\mathbf{y}$ es un
vector de probabilidades que hace justo el valor esperado de todas las apuestas
(condici\'on~1 del teorema). Por el Teorema de Dualidad, si el primal es
factible, el m\'aximo del dual es tambi\'en cero: no es posible una ganancia
garantizada. Si el primal no es factible, el dual no tiene soluci\'on \'optima
finita y la ganancia puede superar cualquier valor positivo
(condici\'on~2 del teorema). $\square$

\begin{intuicion}
La belleza de esta demostraci\'on es que traduce una afirmaci\'on financiera
--- ``existe o no existe arbitraje'' --- en una afirmaci\'on matem\'atica
--- ``el primal es o no es factible''. La dualidad de la programaci\'on lineal
hace el resto. La prueba tambi\'en aclara que las condiciones~1 y 2 del teorema
no son simplemente ``distintas'' sino que una es la negaci\'on exacta de la
otra, en el sentido del \'algebra lineal.
\end{intuicion}

%%% =========================================================
\section{Resumen y Conclusiones}
%%% =========================================================

El cap\'itulo establece el \textbf{Teorema de Arbitraje (Arbitrage Theorem)}
como fundamento matem\'atico de la valoraci\'on libre de arbitraje de
instrumentos financieros. Los puntos centrales son:

\begin{itemize}
  \item El \textbf{arbitraje (arbitrage)} existe cuando los precios de dos
        instrumentos son inconsistentes entre s\'i, permitiendo una ganancia
        sin riesgo. Los mercados eficientes eliminan estas oportunidades al
        atraer inversores hacia los activos con precio incorrecto.

  \item Los \textbf{momios (odds)} $n:m$ se relacionan con probabilidades
        mediante $P = m/(n+m)$. Si las probabilidades impl\'icitas en los
        momios de todos los resultados no suman 1, existe arbitraje. Esta
        incoherencia es detectable algebraicamente a trav\'es de la
        programaci\'on lineal.

  \item La \textbf{programaci\'on lineal (linear programming)} minimiza
        $\mathbf{c} \cdot \mathbf{x}$ sujeto a $A\mathbf{x} = \mathbf{b}$,
        $\mathbf{x} \geq 0$. El conjunto factible es siempre convexo y la
        soluci\'on \'optima se halla en un v\'ertice.

  \item El \textbf{Teorema de Dualidad D\'ebil (Weak Duality Theorem)}
        establece $\mathbf{y} \cdot \mathbf{b} \leq \mathbf{c} \cdot \mathbf{x}$
        para cualquier par de soluciones factibles del primal y el dual.
        La igualdad caracteriza la optimalidad simult\'anea de ambos.

  \item La \textbf{holgura complementaria (complementary slackness)} exige
        que para cada $j$, bien $x_j = 0$ o bien $(\mathbf{y}A)_j = c_j$.
        Esta condici\'on permite determinar las variables nulas del primal
        a partir del \'optimo dual, reduciendo el problema significativamente.

  \item El \textbf{Teorema de Dualidad (Duality Theorem)} garantiza que si
        el primal tiene soluci\'on \'optima entonces el dual tambi\'en la
        tiene, y $\min(\mathbf{c} \cdot \mathbf{x}) = \max(\mathbf{y} \cdot \mathbf{b})$
        --- hay ``brecha cero'' entre ambos problemas.

  \item El \textbf{Teorema de Arbitraje (Arbitrage Theorem)} establece que
        exactamente uno de dos escenarios es posible: (i)~existe un vector de
        probabilidades $\mathbf{y}$ con $\sum_j y_j = 1$ que hace $\sum_j y_j r_{ij} = 0$
        para toda apuesta $i$ --- todos los precios son justos; o (ii)~existe
        una estrategia $\mathbf{x}$ tal que $\sum_i x_i r_{ij} > 0$ para todo
        resultado $j$ --- hay ganancia garantizada.

  \item La condici\'on (i) define la \textbf{medida de probabilidad neutral
        al riesgo (risk-neutral probability measure)}, herramienta central
        de la valoraci\'on de derivados. La ausencia de esta medida implica
        arbitraje. Este principio es la base de la ecuaci\'on de
        \textbf{Black--Scholes} derivada en el Cap\'itulo~6.
\end{itemize}

\end{document}
